\documentclass[10pt]{article}

\usepackage[utf8]{inputenc}

\usepackage[T1]{fontenc}

\usepackage{amsmath}

\usepackage{amsfonts}

\usepackage{amssymb}

\usepackage[version=4]{mhchem}

\usepackage{stmaryrd}

\usepackage{mathrsfs}

\usepackage{bbold}



\begin{document}

пересекающихся петель на $S$. Тогда существует единственная с точностью до конформной әквивалентности стандартная У. $(D, \pi, G)$ поверхности $S$ такая, что каждая факторподгруппа $G$ является либо фуксовоӥ, либо әлементарной и накрытие $\pi: D \rightarrow S$ построено по ваименьшей нормальной подгруппе $\pi_{1}(S)$, натянутой на петли $v_{1}, \ldots, v_{n}$.



Теория квазиконформных отображений и Тайхмюллера пространств позволила доказать возможность одновременной У. нескольких римановых поверхностей одной клейновой группой, а также всех римановых поверхностей данного типа (см. [7]).



Jum.: [1] K 1 e i n C h r. F., "Math. Ann", 1883, Bd 21,

S. 141-218; [2] P o i n c a r e H., "Acta math.", 1907, v. 31, p. 1-63, [3] K o e b e P., "Göttinger Nachr.», 1907, S. 177-210; [4] M a s k i t B., "Ann." Math.", 1965, v. 81, 숭, p. 341-55; [6] е г о ж ж, "Amer. J. Math.", 1968 , v. 90, № 3 , p. 718-22, p. 293-312, [7] Б е p c J., "Успехи матем. наук", 1973 , т. 28, в. 4, с. 153-98, [8] К с о в Б. Н., Г у с е в с к и й Н. А., Клейновы группы и униформизация в примерах и задачах, Новосиб., 1981; [9] Н е в а н-

ЈI и н н а Р., Униформизация, пер с нем., М., 1955, [10] Ф $о$ р д Ј. Р., Автоморфные функции, пер. с англ., М.- Л., 1936. Н. А. Гусевский.



УНИФОРМИЗУЮЩИЙ ЭЛЕМЕНТ - элемент $\pi$ дискретно нормированного кольца $A$ с простым идеалом $\mathfrak{p}$ такой, что $\mathfrak{p}=A \pi$. Если $\pi_{1}, \pi_{2}-$ два У. ә. в $A$, то элемент $\pi_{1} \pi_{2}^{-1}$ обратим в $A$. Пусть $R$ - нек-рая система представителей в $A$ элементов поля вычетов $A / \mathfrak{p}$. Тогда любой элемент $a \in A$ однозначно записывается в виде степенного ряда $\sum_{n=0}^{\infty} r_{i} \pi^{i}$, где $r_{i} \in R, \pi-У$. . Если кольцо $A$ полно относительно дискретного нормирования, то любой степенной ряд указанного вида представляет нек-рый элемент $a \in A$.



Если $A$ - локальное кольцо функций простой точки $x$ алгебраич. кривой $X$, то $\pi$ является $\mathscr{У}$. ә. тогда и только тогда, когда $\pi$ имеет нуль первого порядка в точке $x$.



УНИЧТОЖЕНИЯ ОПЕРАТОРЫ - семейство замкнутых линейных операторов $\{a(f), f \in H\}$, где $H$ - некрое гильбертово пространство, действующих в пространстве, построенном над пространством $H$ (т. е. в симметризованной $\Gamma^{s}(H)$ или антисимметризованной $\Gamma^{a}(H)$ тензорной әкспоненте пространства $\left.H\right)$ так, что на векторах $\left(f_{1} \otimes \ldots \otimes f_{n}\right)_{\alpha} \in \Gamma^{\alpha}(H), \quad \alpha=s, a, \quad$ являющихся симметризованными $(\alpha=s)$ или антисимметризованными $(\alpha=a)$ тензорными произведениями последовательностей элементов $f_{1}, \ldots, f_{n} \in H, n=1,2, \ldots$, из $H$, они задаются формулами:



\[

\begin{gathered}

a(f)\left(f_{1} \otimes \ldots \otimes f_{n}\right)_{s}=\sum_{i=1}^{n}\left(f, f_{i}\right)\left(f_{1} \otimes \ldots \otimes f_{i-1} \otimes\right. \\

\left.\otimes f_{i+1} \otimes \ldots \otimes f_{n}\right)_{s} \\

\text { в симметричном } \begin{array}{c}

\text { случае, и } \\

a(f)\left(f_{1} \otimes \ldots \otimes f_{n}\right)_{a}= \\

=\sum_{i=1}^{n}(-1)^{i-1}\left(f, f_{i}\right)\left(f_{1} \otimes \ldots \otimes f_{i-1} \otimes f_{i+1} \otimes \ldots \otimes f_{n}\right)_{a}

\end{array}

\end{gathered}

\]



в антисимметричном случае; вакуумный вектор $\Omega \in \Gamma^{\alpha}(H), \alpha=s, a$ (т. е. единичный вектор в подпространстве констант из $\left.\Gamma^{\alpha}(H)\right)$ переводится операторами $a(f)$ в нуль. Операторы $\left\{a^{*}(f), f \in H\right\}$, сопряженные к операторам $a(f)$, наз. опе р а тор а и и ро ж д ен и я, действуют на векторы $\left(f_{1} \otimes \ldots \otimes f_{n}\right)_{\alpha}, \alpha=s, a, h=$ $=1,2, \ldots$ по формулам



\[

a^{*}(f)\left(f_{1} \otimes \ldots \otimes f_{n}\right)_{\alpha}=\left(f \otimes f_{1} \otimes \ldots \otimes f_{n}\right)_{\alpha}

\]



и



\[

a^{*}(f) \Omega=f .

\]



При этом для каждого $n>0$ шодпространство $\Gamma_{n}^{\alpha}(H)_{\alpha}^{\bigotimes} n$, $\alpha=s, a$ (симметризованная или антпсимметризованная $n$-я тензорная степень $H)$ переводится оператором $\boldsymbol{a}(f)$ в подпространство $\Gamma_{n-1}^{\alpha}(H)$, а оператором $a^{*}(f)-$ в подпространство $\Gamma_{n+1}^{\alpha}(H)$.



При физич. интерпретации пространства Фока $\Gamma^{\alpha}(H)$ $\alpha=s, a$ как пространства состояний системы, состоящей из произвольного (конечного) числа квантовых частиц с пространством $H$ в качестве пространства состояний отдельной частицы, подпространство $\Gamma_{n}^{\alpha}(H)$ соответствует $n$-частичным состояниям системы, т. е. состояниям, в к-рых имеется ровно $n$ частиц, и, таким образом, $n$-тастичные состояния операторы $a(f)$ переводят в $(n-1)$-частичные («уничтожают» частицу), а операторы $a^{*}(f)-$ в $(n+1)$-частичные ("рождают» частицу).



Операторы $a(f)$ и $a^{*}(f)$ образуют неприводимое семейство операторов и удовлетворяют следующей системе п е р е с т а но во чны симметричном случае



$a\left(f_{1}\right) a\left(f_{2}\right)-a\left(f_{2}\right) a\left(f_{1}\right)=a^{*}\left(f_{1}\right) a^{*}\left(f_{2}\right)-a^{*}\left(f_{2}\right) a^{*}\left(f_{1}\right)=0$,



$a^{*}\left(f_{2}\right) a\left(f_{1}\right)-a\left(f_{1}\right) a^{*}\left(f_{2}\right)=-\left(f_{1}, f_{2}\right) E$



(с о о тно шения комм у т ации), в антисимметричном случае



$a\left(f_{1}\right) a\left(f_{2}\right)+a\left(f_{2}\right) a\left(f_{1}\right)=a^{*}\left(f_{1}\right) a^{*}\left(f_{2}\right)+a^{*}\left(f_{2}\right) a^{*}\left(f_{1}\right)=0$,



\[

a\left(f_{1}\right) a^{*}\left(f_{2}\right)+a^{*}\left(f_{2}\right) a\left(f_{1}\right)=\left(f_{1}, f_{2}\right) E

\]



(с о о тношения антиком м у т а ци и), здесь $(\cdot, \cdot)$ - скалярное произведение в $H$, а $E$ - единичный оператор в $\Gamma^{s}(H)$ или $\Gamma^{a}(H)$. Кроме описанных здесь семейств операторов $\boldsymbol{a}(f)$ и $a^{*}(f), f \in H$ в случае бесконечномерного пространства $H$ существуют и другие, не әквивалентные им, неприводимые представления коммутационных и антикоммутационных соотношений (4) или (5). Иногда их также называют операторами рождения и уничтожения. В случае же конечномерного $H$ все неприводимые представления как коммутационных, так и антикоммутационных соотношений унитарно эквивалентны.



Операторы $\left\{a(f), a^{*}(f), f \in H\right\}$ являются во многих отношениях удобными «образующими» в совокупности всех линейных операторов, действующих в пространстве $\Gamma^{\alpha}(H), \alpha=s, a$, и представление таких операторов в виде сумм произведений огераторов рождения и уничтожения (нормальная форма оператора) бывает часто полезным при их исследовании. Связанный с этим формализм носит название м е т о д а в т 0 p п ч н о г 0 к в а н т о в а ни я, см. (1).



В частном, но важном для применений случае, когда $H=L_{2}\left(\mathbb{R} v, d^{v} x\right), v=1,2, \ldots$, (или в более общем случае $H=L_{2}(M, \sigma)$, где $(M, \sigma)$ - пространство с мерой) семейство операторов $\boldsymbol{a}(f), \quad \boldsymbol{a}^{*}(f), f \in L_{2}\left(\mathbb{R} v, d^{v} x\right)$, порождает две операторнозначные обобщенные функции $a(x)$ и $a^{*}(x)$ такие, что



\[

a(f)=\int_{\mathbb{R}} v^{a}(x) f(x) d^{v} x, a^{*}(f)=\int_{\mathbb{R}} a^{*}(x) \bar{f}(x) d^{v} x .

\]



Введение функций $a(x)$ и $a^{*}(x)$ для формализма вторичного квантования оказывается удобным (напр., позволяет непосредственно рассматривать операторы вида



$\int\left(\mathbb{R}^{v}\right)^{m+n} K\left(x_{1}, \ldots, x_{n} ; y_{1}, \ldots, y_{m}\right) a^{*}\left(x_{1}\right) \ldots a^{*}\left(x_{n}\right) \times$ $\times a\left(y_{1}\right) \ldots a\left(y_{n}\right) d^{v} x_{1} \ldots d^{v} x_{n} d^{v} y_{1} \ldots d^{v} y_{m}, n, m=1,2, \ldots$, где $K\left(x_{1}, \ldots, x_{n} ; y_{1}, \ldots, y_{m}\right)$ - нек-рая «достаточно хорошая》 функция), не прибегая к их разложению в ряды Іо мономам



\[

a^{*}\left(f_{1}\right) \ldots a^{*}\left(f_{n}\right) a\left(g_{1}\right) \ldots a\left(g_{m}\right)

\]



где



$f_{1}, \ldots, f_{n}, g_{1}, \ldots, g_{m} \in L_{2}\left(\mathbb{R} v, d^{v} x\right)$.





\end{document}